\documentclass[12pt]{article}
\usepackage[margin=1in]{geometry}
\usepackage{enumitem}
\usepackage{titlesec}
\usepackage{parskip}
\usepackage{amsmath}
\usepackage{graphicx}
\usepackage{hyperref}
\usepackage{fontenc}

\title{\textbf{Nalepa 2010 RG Notes}\\
\large Nalepa. \textit{Skeletons in the Closet: Transitional Justice in Post-Communist Europe}.\\
New York: Cambridge University Press, 2010.}
\author{}
\date{}

\begin{document}

\maketitle

\section{Main Idea}



%--------------------------------------------------
\section{General Notes}
%--------------------------------------------------

\begin{itemize}[leftmargin=*, itemsep=4pt]
    \item Differentiating why some post-USSR democracies pursued lustration, the discovery/vetting/exposure of those who were collaborators with the secret polise, while others did not (or only did to varying extents)
    \item Game theoretic conclusion -- lustration depends on politicians' private incentives to oppose lustration on the grounds of their own previous transgressions (i.e., the skeletons in the closet)
    \item Those without skeletons stand to gain normatively and politically from pursuign lustration, including in future political races or discussions
    \item Previous literature suggests this should just be in line with a party's previous support for communism; work shows that this is not the case
\end{itemize}


\section{Important Specific Factual Claims}

\begin{itemize}[leftmargin=*, itemsep=4pt]
    \item Motivating instance is the bargaining over transition of power, primarily as it stands to significantly influence whether or not documents/materials are kept or destroyed
\end{itemize}


\section{Connections}

\begin{itemize}[leftmargin=*, itemsep=4pt]
    \item
\end{itemize}


\section{My Critiques}

\begin{itemize}[leftmargin=*, itemsep=4pt]
    \item
\end{itemize}


\end{document}
